% GENERATED FILE. Edit canonical lessons, then run npm run generate:books.
% canonical-source-hash: fnv1a64:633a694aa20090dd
% canonical-lessons: BN-C93-bosonto, BN-C93-hemonto, BN-C93-sekend, BN-C93-shuru, BN-C93-ajkal

\chapter{Spring, Late autumn, Second, Beginning, Nowadays}
\label{ch:bn-spring-late-autumn-second-beginning-nowadays}
\hlchaptermodality{93}

\begin{chapteropening}
\textbf{By the end of this chapter:} \emph{I can say spring, late autumn, a second, the beginning and nowadays.}

Complete the last lesson of chapter 93: I can say spring, late autumn, a second, the beginning and nowadays.
\end{chapteropening}

\section[\texorpdfstring{\bn{বসন্ত} (bôsôntô)}{বসন্ত (bôsôntô)}]{\bn{বসন্ত} (bôsôntô) — spring}

\label{lesson:BN-C93-bosonto}



\begin{quote}
Before the new one: say the Bengali for daily, then the Bengali for winter; cold.
\end{quote}

\subsection*{You'll want to know: \bn{বসন্ত}}
\textbf{\bn{বসন্ত}} — \emph{bôsôntô} — \textquotedblleft{}spring\textquotedblright{}.

\begin{cousinweb}
\textbf{\bn{বসন্ত}} (\emph{bôsôntô}) is Sanskrit \emph{vasanta}, \textquotedblleft{}spring season\textquotedblright{}, taken over whole as a learned loan, a tatsama.
\end{cousinweb}

\begin{grammarlens}[title={What you've built}]
One more word for this chapter.
\end{grammarlens}

\subsection*{Your turn}
\begin{itemize}
\raggedright
  \item \emph{Say it:} \emph{bôsôntô}
  \item \emph{Say it:} \emph{bôsôntô}, once more
  \item \emph{Say it:} say \emph{bôsôntô}
  \item \emph{Recall:} say the Bengali for daily, then the Bengali for winter; cold, then say \emph{bôsôntô} again
\end{itemize}

\begin{tcolorbox}[breakable,colback=teal!4,colframe=teal!35!black,title={Before you move on}]
What does \bn{বসন্ত} mean? (\textquotedblleft{}Spring\textquotedblright{}.) Say it once more.
\end{tcolorbox}
\section[\texorpdfstring{\bn{হেমন্ত} (hemôntô)}{হেমন্ত (hemôntô)}]{\bn{হেমন্ত} (hemôntô) — late autumn}

\label{lesson:BN-C93-hemonto}



\begin{quote}
Before the new one: say the Bengali for winter; cold, then the Bengali for spring.
\end{quote}

\subsection*{You'll want to know: \bn{হেমন্ত}}
\textbf{\bn{হেমন্ত}} — \emph{hemôntô} — \textquotedblleft{}late autumn\textquotedblright{}.

\begin{cousinweb}
\textbf{\bn{হেমন্ত}} (\emph{hemôntô}) is Sanskrit \emph{hemanta}, \textquotedblleft{}winter\textquotedblright{}, taken over whole as a learned loan, a tatsama (CDIAL 14164). The word inherited from it is \emph{heõt}, \textquotedblleft{}winter\textquotedblright{}, which CDIAL lists; in Bengali the season it names is late autumn.
\end{cousinweb}

\begin{grammarlens}[title={What you've built}]
One more word for this chapter.
\end{grammarlens}

\subsection*{Your turn}
\begin{itemize}
\raggedright
  \item \emph{Say it:} \emph{hemôntô}
  \item \emph{Say it:} \emph{hemôntô}, once more
  \item \emph{Say it:} say \emph{hemôntô}
  \item \emph{Recall:} say the Bengali for winter; cold, then the Bengali for spring, then say \emph{hemôntô} again
\end{itemize}

\begin{tcolorbox}[breakable,colback=teal!4,colframe=teal!35!black,title={Before you move on}]
What does \bn{হেমন্ত} mean? (\textquotedblleft{}Late autumn\textquotedblright{}.) Say it once more.
\end{tcolorbox}
\section[\texorpdfstring{\bn{সেকেন্ড} (sekend)}{সেকেন্ড (sekend)}]{\bn{সেকেন্ড} (sekend) — a second}

\label{lesson:BN-C93-sekend}



\begin{quote}
Before the new one: say the Bengali for spring, then the Bengali for late autumn.
\end{quote}

\subsection*{You'll want to know: \bn{সেকেন্ড}}
\textbf{\bn{সেকেন্ড}} — \emph{sekend} — \textquotedblleft{}a second\textquotedblright{}.

\begin{cousinweb}
\textbf{\bn{সেকেন্ড}} (\emph{sekend}) is a loan from English \emph{second}.
\end{cousinweb}

\begin{grammarlens}[title={What you've built}]
One more word for this chapter.
\end{grammarlens}

\subsection*{Your turn}
\begin{itemize}
\raggedright
  \item \emph{Say it:} \emph{sekend}
  \item \emph{Say it:} \emph{sekend}, once more
  \item \emph{Say it:} say \emph{sekend}
  \item \emph{Recall:} say the Bengali for spring, then the Bengali for late autumn, then say \emph{sekend} again
\end{itemize}

\begin{tcolorbox}[breakable,colback=teal!4,colframe=teal!35!black,title={Before you move on}]
What does \bn{সেকেন্ড} mean? (\textquotedblleft{}A second\textquotedblright{}.) Say it once more.
\end{tcolorbox}
\section[\texorpdfstring{\bn{শুরু} (shuru)}{শুরু (shuru)}]{\bn{শুরু} (shuru) — the beginning}

\label{lesson:BN-C93-shuru}



\begin{quote}
Before the new one: say the Bengali for late autumn, then the Bengali for a second.
\end{quote}

\subsection*{You'll want to know: \bn{শুরু}}
\textbf{\bn{শুরু}} — \emph{shuru} — \textquotedblleft{}the beginning\textquotedblright{}.

\begin{cousinweb}
\textbf{\bn{শুরু}} (\emph{shuru}) is a loan from Arabic \emph{shurūʿ}, \textquotedblleft{}beginning\textquotedblright{}, which reached Bengali through Persian. Its Arabic root is \emph{sh-r-ʿ} (Lane).
\end{cousinweb}

\begin{grammarlens}[title={What you've built}]
One more word for this chapter.
\end{grammarlens}

\subsection*{Your turn}
\begin{itemize}
\raggedright
  \item \emph{Say it:} \emph{shuru}
  \item \emph{Say it:} \emph{shuru}, once more
  \item \emph{Say it:} say \emph{shuru}
  \item \emph{Recall:} say the Bengali for late autumn, then the Bengali for a second, then say \emph{shuru} again
\end{itemize}

\begin{tcolorbox}[breakable,colback=teal!4,colframe=teal!35!black,title={Before you move on}]
What does \bn{শুরু} mean? (\textquotedblleft{}The beginning\textquotedblright{}.) Say it once more.
\end{tcolorbox}
\section[\texorpdfstring{\bn{আজকাল} (ājkāl)}{আজকাল (ājkāl)}]{\bn{আজকাল} (ājkāl) — nowadays}

\label{lesson:BN-C93-ajkal}



\begin{quote}
Before the new one: say the Bengali for a second, then the Bengali for the beginning.
\end{quote}

\subsection*{You'll want to know: \bn{আজকাল}}
\textbf{\bn{আজকাল}} — \emph{ājkāl} — \textquotedblleft{}nowadays\textquotedblright{}.

\begin{cousinweb}
\textbf{\bn{আজকাল}} (\emph{ājkāl}) is \emph{āj}, \textquotedblleft{}today\textquotedblright{}, from Sanskrit \emph{adya} (CDIAL 242), + \emph{kāl}, \textquotedblleft{}tomorrow, yesterday\textquotedblright{}, from Sanskrit \emph{kalya}, \textquotedblleft{}dawn\textquotedblright{} (CDIAL 3104): \textquotedblleft{}today and tomorrow\textquotedblright{}, so \textquotedblleft{}these days\textquotedblright{}.
\end{cousinweb}

\begin{grammarlens}[title={What you've built}]
One more word for this chapter.
\end{grammarlens}

\subsection*{Your turn}
\begin{itemize}
\raggedright
  \item \emph{Say it:} \emph{ājkāl}
  \item \emph{Say it:} \emph{ājkāl}, once more
  \item \emph{Say it:} say \emph{ājkāl}
  \item \emph{Recall:} say the Bengali for a second, then the Bengali for the beginning, then say \emph{ājkāl} again
\end{itemize}

\begin{tcolorbox}[breakable,colback=teal!4,colframe=teal!35!black,title={Before you move on}]
What does \bn{আজকাল} mean? (\textquotedblleft{}Nowadays\textquotedblright{}.) Say it once more.
\end{tcolorbox}
